%% ECON 282E -- Session 3, Deck A2: Value Function Iteration on a Grid
%% Budget: 16 frames (~22 minutes). Provenance: slides/SOURCES.md.
%% Every number and every plotted coordinate on this deck comes from
%% tools/figures/s03_numbers.json, generated by tools/figures/s03_generate.py on a
%% compute node (2026-09-17).  Calibration: u=ln c, y=z k^alpha, delta=1,
%% alpha=0.36, beta=0.95, k in [0.05,0.60], N=200, tol=1e-6, float64.

%% ECON 282E -- Foundations of Macroeconomics (UC Riverside, Fall 2026)
%% Shared Beamer preamble for every deck in slides/S01 ... slides/S10.
%%
%% Usage (a deck lives in slides/S0X/ and is compiled from that folder):
%%   %% ECON 282E -- Foundations of Macroeconomics (UC Riverside, Fall 2026)
%% Shared Beamer preamble for every deck in slides/S01 ... slides/S10.
%%
%% Usage (a deck lives in slides/S0X/ and is compiled from that folder):
%%   %% ECON 282E -- Foundations of Macroeconomics (UC Riverside, Fall 2026)
%% Shared Beamer preamble for every deck in slides/S01 ... slides/S10.
%%
%% Usage (a deck lives in slides/S0X/ and is compiled from that folder):
%%   \input{../preamble/course_preamble.tex}
%%   \decktitle{1}{A1}{Who I Am \& Why This Course}{September 24, 2026}
%%   \begin{document}
%%   \makedecktitle
%%   ...
%%
%% Adapted from Courses/AI-ECON-2026/source/shared_preamble.tex (Metropolis theme,
%% concept/caution/example/takeaway boxes, \sectiondivider). Changes for this course:
%%   * course title block (\decktitle / \makedecktitle) instead of \makebeamertitle
%%   * \zh{...} typesets Chinese (book titles) under pdflatex via CJKutf8
%%   * researchbox for the "Research in the wild" frames
%%   * section pages off; TikZ libraries and the course map loaded here
%% Build with pdflatex only (two passes); tools/build_slides.py does this for you.

\documentclass[10pt,english,aspectratio=169]{beamer}

\usepackage[T1]{fontenc}
\usepackage{textcomp}
\usepackage[utf8]{inputenc}
\usepackage{babel}
\usepackage{CJKutf8}

\usepackage{amsmath, amssymb, amsthm, graphicx}
\usepackage{booktabs, multirow, tabularx}
\usepackage{tikz}
\usetikzlibrary{positioning,arrows.meta,shapes,calc,fit,backgrounds}
\usepackage{pgfplots}
\pgfplotsset{compat=1.16}
\usepackage[most]{tcolorbox}
\tcbuselibrary{skins, breakable}
\usepackage{listings}
\usepackage{xcolor}

\lstset{
  basicstyle=\ttfamily\small,
  keywordstyle=\color{blue!80!black}\bfseries,
  commentstyle=\color{green!50!black}\itshape,
  stringstyle=\color{orange!80!black},
  backgroundcolor=\color{gray!8},
  frame=single,
  rulecolor=\color{gray!40},
  breaklines=true,
  showstringspaces=false,
  numbers=none,
}

\usetheme{metropolis}
\usefonttheme{professionalfonts}
\metroset{sectionpage=none}

\definecolor{LightGray}{RGB}{230,230,230}
\definecolor{RedTitle}{RGB}{0,0,200}
\definecolor{VeryLightGray}{RGB}{245,245,245}
\definecolor{DarkText}{RGB}{0,0,0}
\definecolor{UCRBlue}{RGB}{0,45,106}
\definecolor{UCRGold}{RGB}{241,171,0}

% Stage colours of the course arc (used by the course map and elsewhere)
\colorlet{StageTrunk}{blue!12}
\colorlet{StageBranch}{cyan!14}
\colorlet{StageMachine}{gray!18}
\colorlet{StageWall}{red!14}
\colorlet{StageTools}{orange!20}
\colorlet{StageData}{green!16}
\colorlet{StageAgents}{violet!14}

\hypersetup{bookmarks=false,breaklinks=true,pdfborder={0 0 0},colorlinks=true,
  linkcolor=DarkText,urlcolor=RedTitle}

\setbeamercolor{background canvas}{bg=white}
\setbeamercolor{title}{fg=RedTitle}
\setbeamercolor{frametitle}{bg=VeryLightGray,fg=DarkText}
\setbeamercolor{normal text}{fg=DarkText}
\setbeamercolor{alerted text}{fg=RedTitle}
\setbeamersize{text margin left=5mm, text margin right=5mm}
\addtobeamertemplate{frametitle}{}{\vspace{-0.5em}}

\setbeamertemplate{navigation symbols}{}
\setbeamertemplate{footline}{%
  \hspace{5mm}{\usebeamerfont{page number in head/foot}\color{black!45}%
  ECON 282E \,$\cdot$\, \insertshortdeck}\hfill%
  \usebeamercolor[fg]{page number in head/foot}%
  \usebeamerfont{page number in head/foot}%
  \insertframenumber\,/\,\inserttotalframenumber\kern1em\vskip2pt%
}

% ---------------------------------------------------------------------------
% Chinese text under pdflatex (book titles etc.)
\newcommand{\zh}[1]{\begin{CJK*}{UTF8}{gbsn}#1\end{CJK*}}

% ---------------------------------------------------------------------------
% Deck title block
%   \decktitle{session}{deck id}{title}{date}
\newcommand{\insertshortdeck}{}
\newcommand{\decksession}{}
\newcommand{\deckid}{}
\newcommand{\decktitletext}{}
\newcommand{\deckdate}{}
\newcommand{\decktitle}[4]{%
  \renewcommand{\decksession}{#1}%
  \renewcommand{\deckid}{#2}%
  \renewcommand{\decktitletext}{#3}%
  \renewcommand{\deckdate}{#4}%
  \renewcommand{\insertshortdeck}{Session #1 \,$\cdot$\, Deck #1.#2}%
  \title{#3}%
  \author{Zhigang Feng}%
  \date{#4}%
}
\newcommand{\makedecktitle}{%
  \begin{frame}[plain,noframenumbering]
    \begin{tikzpicture}[remember picture,overlay]
      \fill[UCRBlue] (current page.north west) rectangle ([yshift=-0.55cm]current page.north east);
      \fill[UCRGold] ([yshift=-0.55cm]current page.north west) rectangle ([yshift=-0.62cm]current page.north east);
      \node[anchor=west, text=white, font=\small\bfseries] at ([xshift=5mm,yshift=-0.275cm]current page.north west)
        {ECON 282E \,$\cdot$\, Foundations of Macroeconomics};
      \node[anchor=east, text=white, font=\small] at ([xshift=-5mm,yshift=-0.275cm]current page.north east)
        {UC Riverside \,$\cdot$\, Fall 2026};
    \end{tikzpicture}
    \vspace*{1.1cm}
    {\usebeamerfont{title}\color{black!55}\large Session \decksession\ \,$\cdot$\, Deck \decksession.\deckid\par}
    \vspace{0.25cm}
    {\usebeamerfont{title}\color{RedTitle}\LARGE\bfseries \decktitletext\par}
    \vspace{0.7cm}
    {\large Zhigang Feng\par}
    \vspace{0.15cm}
    {\color{black!65}\deckdate\par}
    \vfill
    {\footnotesize\itshape\color{black!55}Build it on paper $\;\to\;$ solve it on the machine $\;\to\;$ push it to the frontier with AI\par}
  \end{frame}}

% ---------------------------------------------------------------------------
% Section divider (plain frame, not counted)
\newcommand{\sectiondivider}[2]{%
\begin{frame}[plain,noframenumbering]
\begin{center}
\vfill
{\Large\textbf{#1}\par}
\vspace{0.35cm}
{\normalsize #2\par}
\vfill
\end{center}
\end{frame}
}

% ---------------------------------------------------------------------------
% Boxes
\newtcolorbox{conceptbox}[1][]{colback=blue!3, colframe=RedTitle!55, boxrule=0.35mm,
  arc=1mm, left=2mm, right=2mm, top=1mm, bottom=1mm, #1}
\newtcolorbox{cautionbox}[1][]{colback=red!3, colframe=red!50!black, boxrule=0.35mm,
  arc=1mm, left=2mm, right=2mm, top=1mm, bottom=1mm, #1}
\newtcolorbox{examplebox}[1][]{colback=green!3, colframe=green!45!black, boxrule=0.35mm,
  arc=1mm, left=2mm, right=2mm, top=1mm, bottom=1mm, #1}
\newtcolorbox{takeawaybox}[1][]{colback=VeryLightGray, colframe=RedTitle!55, boxrule=0.35mm,
  arc=1mm, left=2mm, right=2mm, top=1mm, bottom=1mm, #1}
\newtcolorbox{researchbox}[1][]{enhanced, colback=violet!4, colframe=violet!60!black,
  boxrule=0.35mm, arc=1mm, left=2mm, right=2mm, top=1mm, bottom=1mm,
  fonttitle=\bfseries\small, title={Research in the wild}, #1}

\newcommand{\compacttables}{\renewcommand{\arraystretch}{1.15}}
\newcommand{\src}[1]{{\tiny\color{black!55}#1}}

% ---------------------------------------------------------------------------
\input{../preamble/course_map.tex}

%%   \decktitle{1}{A1}{Who I Am \& Why This Course}{September 24, 2026}
%%   \begin{document}
%%   \makedecktitle
%%   ...
%%
%% Adapted from Courses/AI-ECON-2026/source/shared_preamble.tex (Metropolis theme,
%% concept/caution/example/takeaway boxes, \sectiondivider). Changes for this course:
%%   * course title block (\decktitle / \makedecktitle) instead of \makebeamertitle
%%   * \zh{...} typesets Chinese (book titles) under pdflatex via CJKutf8
%%   * researchbox for the "Research in the wild" frames
%%   * section pages off; TikZ libraries and the course map loaded here
%% Build with pdflatex only (two passes); tools/build_slides.py does this for you.

\documentclass[10pt,english,aspectratio=169]{beamer}

\usepackage[T1]{fontenc}
\usepackage{textcomp}
\usepackage[utf8]{inputenc}
\usepackage{babel}
\usepackage{CJKutf8}

\usepackage{amsmath, amssymb, amsthm, graphicx}
\usepackage{booktabs, multirow, tabularx}
\usepackage{tikz}
\usetikzlibrary{positioning,arrows.meta,shapes,calc,fit,backgrounds}
\usepackage{pgfplots}
\pgfplotsset{compat=1.16}
\usepackage[most]{tcolorbox}
\tcbuselibrary{skins, breakable}
\usepackage{listings}
\usepackage{xcolor}

\lstset{
  basicstyle=\ttfamily\small,
  keywordstyle=\color{blue!80!black}\bfseries,
  commentstyle=\color{green!50!black}\itshape,
  stringstyle=\color{orange!80!black},
  backgroundcolor=\color{gray!8},
  frame=single,
  rulecolor=\color{gray!40},
  breaklines=true,
  showstringspaces=false,
  numbers=none,
}

\usetheme{metropolis}
\usefonttheme{professionalfonts}
\metroset{sectionpage=none}

\definecolor{LightGray}{RGB}{230,230,230}
\definecolor{RedTitle}{RGB}{0,0,200}
\definecolor{VeryLightGray}{RGB}{245,245,245}
\definecolor{DarkText}{RGB}{0,0,0}
\definecolor{UCRBlue}{RGB}{0,45,106}
\definecolor{UCRGold}{RGB}{241,171,0}

% Stage colours of the course arc (used by the course map and elsewhere)
\colorlet{StageTrunk}{blue!12}
\colorlet{StageBranch}{cyan!14}
\colorlet{StageMachine}{gray!18}
\colorlet{StageWall}{red!14}
\colorlet{StageTools}{orange!20}
\colorlet{StageData}{green!16}
\colorlet{StageAgents}{violet!14}

\hypersetup{bookmarks=false,breaklinks=true,pdfborder={0 0 0},colorlinks=true,
  linkcolor=DarkText,urlcolor=RedTitle}

\setbeamercolor{background canvas}{bg=white}
\setbeamercolor{title}{fg=RedTitle}
\setbeamercolor{frametitle}{bg=VeryLightGray,fg=DarkText}
\setbeamercolor{normal text}{fg=DarkText}
\setbeamercolor{alerted text}{fg=RedTitle}
\setbeamersize{text margin left=5mm, text margin right=5mm}
\addtobeamertemplate{frametitle}{}{\vspace{-0.5em}}

\setbeamertemplate{navigation symbols}{}
\setbeamertemplate{footline}{%
  \hspace{5mm}{\usebeamerfont{page number in head/foot}\color{black!45}%
  ECON 282E \,$\cdot$\, \insertshortdeck}\hfill%
  \usebeamercolor[fg]{page number in head/foot}%
  \usebeamerfont{page number in head/foot}%
  \insertframenumber\,/\,\inserttotalframenumber\kern1em\vskip2pt%
}

% ---------------------------------------------------------------------------
% Chinese text under pdflatex (book titles etc.)
\newcommand{\zh}[1]{\begin{CJK*}{UTF8}{gbsn}#1\end{CJK*}}

% ---------------------------------------------------------------------------
% Deck title block
%   \decktitle{session}{deck id}{title}{date}
\newcommand{\insertshortdeck}{}
\newcommand{\decksession}{}
\newcommand{\deckid}{}
\newcommand{\decktitletext}{}
\newcommand{\deckdate}{}
\newcommand{\decktitle}[4]{%
  \renewcommand{\decksession}{#1}%
  \renewcommand{\deckid}{#2}%
  \renewcommand{\decktitletext}{#3}%
  \renewcommand{\deckdate}{#4}%
  \renewcommand{\insertshortdeck}{Session #1 \,$\cdot$\, Deck #1.#2}%
  \title{#3}%
  \author{Zhigang Feng}%
  \date{#4}%
}
\newcommand{\makedecktitle}{%
  \begin{frame}[plain,noframenumbering]
    \begin{tikzpicture}[remember picture,overlay]
      \fill[UCRBlue] (current page.north west) rectangle ([yshift=-0.55cm]current page.north east);
      \fill[UCRGold] ([yshift=-0.55cm]current page.north west) rectangle ([yshift=-0.62cm]current page.north east);
      \node[anchor=west, text=white, font=\small\bfseries] at ([xshift=5mm,yshift=-0.275cm]current page.north west)
        {ECON 282E \,$\cdot$\, Foundations of Macroeconomics};
      \node[anchor=east, text=white, font=\small] at ([xshift=-5mm,yshift=-0.275cm]current page.north east)
        {UC Riverside \,$\cdot$\, Fall 2026};
    \end{tikzpicture}
    \vspace*{1.1cm}
    {\usebeamerfont{title}\color{black!55}\large Session \decksession\ \,$\cdot$\, Deck \decksession.\deckid\par}
    \vspace{0.25cm}
    {\usebeamerfont{title}\color{RedTitle}\LARGE\bfseries \decktitletext\par}
    \vspace{0.7cm}
    {\large Zhigang Feng\par}
    \vspace{0.15cm}
    {\color{black!65}\deckdate\par}
    \vfill
    {\footnotesize\itshape\color{black!55}Build it on paper $\;\to\;$ solve it on the machine $\;\to\;$ push it to the frontier with AI\par}
  \end{frame}}

% ---------------------------------------------------------------------------
% Section divider (plain frame, not counted)
\newcommand{\sectiondivider}[2]{%
\begin{frame}[plain,noframenumbering]
\begin{center}
\vfill
{\Large\textbf{#1}\par}
\vspace{0.35cm}
{\normalsize #2\par}
\vfill
\end{center}
\end{frame}
}

% ---------------------------------------------------------------------------
% Boxes
\newtcolorbox{conceptbox}[1][]{colback=blue!3, colframe=RedTitle!55, boxrule=0.35mm,
  arc=1mm, left=2mm, right=2mm, top=1mm, bottom=1mm, #1}
\newtcolorbox{cautionbox}[1][]{colback=red!3, colframe=red!50!black, boxrule=0.35mm,
  arc=1mm, left=2mm, right=2mm, top=1mm, bottom=1mm, #1}
\newtcolorbox{examplebox}[1][]{colback=green!3, colframe=green!45!black, boxrule=0.35mm,
  arc=1mm, left=2mm, right=2mm, top=1mm, bottom=1mm, #1}
\newtcolorbox{takeawaybox}[1][]{colback=VeryLightGray, colframe=RedTitle!55, boxrule=0.35mm,
  arc=1mm, left=2mm, right=2mm, top=1mm, bottom=1mm, #1}
\newtcolorbox{researchbox}[1][]{enhanced, colback=violet!4, colframe=violet!60!black,
  boxrule=0.35mm, arc=1mm, left=2mm, right=2mm, top=1mm, bottom=1mm,
  fonttitle=\bfseries\small, title={Research in the wild}, #1}

\newcommand{\compacttables}{\renewcommand{\arraystretch}{1.15}}
\newcommand{\src}[1]{{\tiny\color{black!55}#1}}

% ---------------------------------------------------------------------------
%% Course map: the recurring "you are here" slide for ECON 282E.
%% Loaded by course_preamble.tex. Pattern follows AI-ECON-2026 lec01_map.tex, but the map
%% shows the whole ten-session course rather than one lecture.
%%
%%   \CourseMap{s}{label}   s = 1..10 highlights session s and prints "you are here: label"
%%                          (e.g. \CourseMap{1}{1.A2}); earlier sessions get a check mark.
%%   \CourseMap{0}{}        full overview, nothing highlighted.
%%
%% Note: TikZ option lists are comma-split before TeX conditionals run, so every \ifnum
%% below sits outside [...] option lists.

\newcommand{\CourseMap}[2]{%
\begin{center}
\begin{tikzpicture}[
    sess/.style={rectangle, rounded corners=2pt, draw=black!45, minimum width=1.34cm,
        minimum height=1.12cm, text width=1.26cm, align=center, inner sep=1pt},
    stage/.style={font=\tiny\bfseries, text=black!70},
]
  \foreach \i/\d/\t/\c in {%
      1/Sep 24/Intro \\ Growth/StageTrunk,
      2/Oct 1/HA $\cdot$ OLG \\ Policy/StageBranch,
      3/Oct 8/Num.\ DP \\ Python/StageMachine,
      4/Oct 15/Numerics \\ I/StageMachine,
      5/Oct 22/Numerics \\ II/StageMachine,
      6/Oct 29/The wall \\ \textbf{Midterm}/StageWall,
      7/Nov 5/ML \\ Deep nets/StageTools,
      8/Nov 12/RL \\ HA $+$ DL/StageTools,
      9/Nov 19/Text \\ LLMs/StageData,
      10/Dec 3/Agents \\ \textbf{Projects}/StageAgents} {
    \node[sess, fill=\c] (n\i) at ({1.46*(\i-1)}, 0)
      {{\scriptsize\bfseries S\i}\\[-1pt]{\tiny\color{black!60}\d}\\[1pt]{\tiny \t}};
    \ifnum\i<#1
      \node[font=\scriptsize, text=green!45!black] at ([xshift=-2.5mm,yshift=-2.2mm]n\i.north east) {$\checkmark$};
    \fi
    \ifnum\i=#1
      \draw[RedTitle, very thick, rounded corners=2pt]
        ([xshift=-0.6mm,yshift=-0.6mm]n\i.south west) rectangle ([xshift=0.6mm,yshift=0.6mm]n\i.north east);
      % keep the label inside the picture at both ends of the row
      \ifnum\i<3
        \node[font=\scriptsize\bfseries, text=RedTitle, anchor=south west] at ([yshift=1.2mm]n\i.north west)
          {$\blacktriangledown$\,you are here: #2};
      \else\ifnum\i>8
        \node[font=\scriptsize\bfseries, text=RedTitle, anchor=south east] at ([yshift=1.2mm]n\i.north east)
          {you are here: #2\,$\blacktriangledown$};
      \else
        \node[font=\scriptsize\bfseries, text=RedTitle, anchor=south] at ([yshift=1.2mm]n\i.north)
          {you are here: #2\,$\blacktriangledown$};
      \fi\fi
    \fi
  }
  % stage band under the sessions
  \foreach \a/\b/\lab/\c in {1/1/trunk/StageTrunk, 2/2/branches/StageBranch,
      3/5/the machine $\cdot$ classical methods/StageMachine, 6/6/the wall/StageWall,
      7/8/new tools/StageTools, 9/9/new data/StageData, 10/10/colleagues/StageAgents} {
    \fill[\c] ([yshift=-1.5mm]n\a.south west) rectangle ([yshift=-4.6mm]n\b.south east);
    \path ([yshift=-1.5mm]n\a.south west) -- ([yshift=-4.6mm]n\b.south east)
      node[midway, stage] {\lab};
  }
  \node[font=\footnotesize\itshape, text=black!70] at ([yshift=-9.5mm]$(n1.south)!0.5!(n10.south)$)
    {Build it on paper $\;\to\;$ solve it on the machine $\;\to\;$ push it to the frontier with AI};
\end{tikzpicture}
\end{center}
}


%%   \decktitle{1}{A1}{Who I Am \& Why This Course}{September 24, 2026}
%%   \begin{document}
%%   \makedecktitle
%%   ...
%%
%% Adapted from Courses/AI-ECON-2026/source/shared_preamble.tex (Metropolis theme,
%% concept/caution/example/takeaway boxes, \sectiondivider). Changes for this course:
%%   * course title block (\decktitle / \makedecktitle) instead of \makebeamertitle
%%   * \zh{...} typesets Chinese (book titles) under pdflatex via CJKutf8
%%   * researchbox for the "Research in the wild" frames
%%   * section pages off; TikZ libraries and the course map loaded here
%% Build with pdflatex only (two passes); tools/build_slides.py does this for you.

\documentclass[10pt,english,aspectratio=169]{beamer}

\usepackage[T1]{fontenc}
\usepackage{textcomp}
\usepackage[utf8]{inputenc}
\usepackage{babel}
\usepackage{CJKutf8}

\usepackage{amsmath, amssymb, amsthm, graphicx}
\usepackage{booktabs, multirow, tabularx}
\usepackage{tikz}
\usetikzlibrary{positioning,arrows.meta,shapes,calc,fit,backgrounds}
\usepackage{pgfplots}
\pgfplotsset{compat=1.16}
\usepackage[most]{tcolorbox}
\tcbuselibrary{skins, breakable}
\usepackage{listings}
\usepackage{xcolor}

\lstset{
  basicstyle=\ttfamily\small,
  keywordstyle=\color{blue!80!black}\bfseries,
  commentstyle=\color{green!50!black}\itshape,
  stringstyle=\color{orange!80!black},
  backgroundcolor=\color{gray!8},
  frame=single,
  rulecolor=\color{gray!40},
  breaklines=true,
  showstringspaces=false,
  numbers=none,
}

\usetheme{metropolis}
\usefonttheme{professionalfonts}
\metroset{sectionpage=none}

\definecolor{LightGray}{RGB}{230,230,230}
\definecolor{RedTitle}{RGB}{0,0,200}
\definecolor{VeryLightGray}{RGB}{245,245,245}
\definecolor{DarkText}{RGB}{0,0,0}
\definecolor{UCRBlue}{RGB}{0,45,106}
\definecolor{UCRGold}{RGB}{241,171,0}

% Stage colours of the course arc (used by the course map and elsewhere)
\colorlet{StageTrunk}{blue!12}
\colorlet{StageBranch}{cyan!14}
\colorlet{StageMachine}{gray!18}
\colorlet{StageWall}{red!14}
\colorlet{StageTools}{orange!20}
\colorlet{StageData}{green!16}
\colorlet{StageAgents}{violet!14}

\hypersetup{bookmarks=false,breaklinks=true,pdfborder={0 0 0},colorlinks=true,
  linkcolor=DarkText,urlcolor=RedTitle}

\setbeamercolor{background canvas}{bg=white}
\setbeamercolor{title}{fg=RedTitle}
\setbeamercolor{frametitle}{bg=VeryLightGray,fg=DarkText}
\setbeamercolor{normal text}{fg=DarkText}
\setbeamercolor{alerted text}{fg=RedTitle}
\setbeamersize{text margin left=5mm, text margin right=5mm}
\addtobeamertemplate{frametitle}{}{\vspace{-0.5em}}

\setbeamertemplate{navigation symbols}{}
\setbeamertemplate{footline}{%
  \hspace{5mm}{\usebeamerfont{page number in head/foot}\color{black!45}%
  ECON 282E \,$\cdot$\, \insertshortdeck}\hfill%
  \usebeamercolor[fg]{page number in head/foot}%
  \usebeamerfont{page number in head/foot}%
  \insertframenumber\,/\,\inserttotalframenumber\kern1em\vskip2pt%
}

% ---------------------------------------------------------------------------
% Chinese text under pdflatex (book titles etc.)
\newcommand{\zh}[1]{\begin{CJK*}{UTF8}{gbsn}#1\end{CJK*}}

% ---------------------------------------------------------------------------
% Deck title block
%   \decktitle{session}{deck id}{title}{date}
\newcommand{\insertshortdeck}{}
\newcommand{\decksession}{}
\newcommand{\deckid}{}
\newcommand{\decktitletext}{}
\newcommand{\deckdate}{}
\newcommand{\decktitle}[4]{%
  \renewcommand{\decksession}{#1}%
  \renewcommand{\deckid}{#2}%
  \renewcommand{\decktitletext}{#3}%
  \renewcommand{\deckdate}{#4}%
  \renewcommand{\insertshortdeck}{Session #1 \,$\cdot$\, Deck #1.#2}%
  \title{#3}%
  \author{Zhigang Feng}%
  \date{#4}%
}
\newcommand{\makedecktitle}{%
  \begin{frame}[plain,noframenumbering]
    \begin{tikzpicture}[remember picture,overlay]
      \fill[UCRBlue] (current page.north west) rectangle ([yshift=-0.55cm]current page.north east);
      \fill[UCRGold] ([yshift=-0.55cm]current page.north west) rectangle ([yshift=-0.62cm]current page.north east);
      \node[anchor=west, text=white, font=\small\bfseries] at ([xshift=5mm,yshift=-0.275cm]current page.north west)
        {ECON 282E \,$\cdot$\, Foundations of Macroeconomics};
      \node[anchor=east, text=white, font=\small] at ([xshift=-5mm,yshift=-0.275cm]current page.north east)
        {UC Riverside \,$\cdot$\, Fall 2026};
    \end{tikzpicture}
    \vspace*{1.1cm}
    {\usebeamerfont{title}\color{black!55}\large Session \decksession\ \,$\cdot$\, Deck \decksession.\deckid\par}
    \vspace{0.25cm}
    {\usebeamerfont{title}\color{RedTitle}\LARGE\bfseries \decktitletext\par}
    \vspace{0.7cm}
    {\large Zhigang Feng\par}
    \vspace{0.15cm}
    {\color{black!65}\deckdate\par}
    \vfill
    {\footnotesize\itshape\color{black!55}Build it on paper $\;\to\;$ solve it on the machine $\;\to\;$ push it to the frontier with AI\par}
  \end{frame}}

% ---------------------------------------------------------------------------
% Section divider (plain frame, not counted)
\newcommand{\sectiondivider}[2]{%
\begin{frame}[plain,noframenumbering]
\begin{center}
\vfill
{\Large\textbf{#1}\par}
\vspace{0.35cm}
{\normalsize #2\par}
\vfill
\end{center}
\end{frame}
}

% ---------------------------------------------------------------------------
% Boxes
\newtcolorbox{conceptbox}[1][]{colback=blue!3, colframe=RedTitle!55, boxrule=0.35mm,
  arc=1mm, left=2mm, right=2mm, top=1mm, bottom=1mm, #1}
\newtcolorbox{cautionbox}[1][]{colback=red!3, colframe=red!50!black, boxrule=0.35mm,
  arc=1mm, left=2mm, right=2mm, top=1mm, bottom=1mm, #1}
\newtcolorbox{examplebox}[1][]{colback=green!3, colframe=green!45!black, boxrule=0.35mm,
  arc=1mm, left=2mm, right=2mm, top=1mm, bottom=1mm, #1}
\newtcolorbox{takeawaybox}[1][]{colback=VeryLightGray, colframe=RedTitle!55, boxrule=0.35mm,
  arc=1mm, left=2mm, right=2mm, top=1mm, bottom=1mm, #1}
\newtcolorbox{researchbox}[1][]{enhanced, colback=violet!4, colframe=violet!60!black,
  boxrule=0.35mm, arc=1mm, left=2mm, right=2mm, top=1mm, bottom=1mm,
  fonttitle=\bfseries\small, title={Research in the wild}, #1}

\newcommand{\compacttables}{\renewcommand{\arraystretch}{1.15}}
\newcommand{\src}[1]{{\tiny\color{black!55}#1}}

% ---------------------------------------------------------------------------
%% Course map: the recurring "you are here" slide for ECON 282E.
%% Loaded by course_preamble.tex. Pattern follows AI-ECON-2026 lec01_map.tex, but the map
%% shows the whole ten-session course rather than one lecture.
%%
%%   \CourseMap{s}{label}   s = 1..10 highlights session s and prints "you are here: label"
%%                          (e.g. \CourseMap{1}{1.A2}); earlier sessions get a check mark.
%%   \CourseMap{0}{}        full overview, nothing highlighted.
%%
%% Note: TikZ option lists are comma-split before TeX conditionals run, so every \ifnum
%% below sits outside [...] option lists.

\newcommand{\CourseMap}[2]{%
\begin{center}
\begin{tikzpicture}[
    sess/.style={rectangle, rounded corners=2pt, draw=black!45, minimum width=1.34cm,
        minimum height=1.12cm, text width=1.26cm, align=center, inner sep=1pt},
    stage/.style={font=\tiny\bfseries, text=black!70},
]
  \foreach \i/\d/\t/\c in {%
      1/Sep 24/Intro \\ Growth/StageTrunk,
      2/Oct 1/HA $\cdot$ OLG \\ Policy/StageBranch,
      3/Oct 8/Num.\ DP \\ Python/StageMachine,
      4/Oct 15/Numerics \\ I/StageMachine,
      5/Oct 22/Numerics \\ II/StageMachine,
      6/Oct 29/The wall \\ \textbf{Midterm}/StageWall,
      7/Nov 5/ML \\ Deep nets/StageTools,
      8/Nov 12/RL \\ HA $+$ DL/StageTools,
      9/Nov 19/Text \\ LLMs/StageData,
      10/Dec 3/Agents \\ \textbf{Projects}/StageAgents} {
    \node[sess, fill=\c] (n\i) at ({1.46*(\i-1)}, 0)
      {{\scriptsize\bfseries S\i}\\[-1pt]{\tiny\color{black!60}\d}\\[1pt]{\tiny \t}};
    \ifnum\i<#1
      \node[font=\scriptsize, text=green!45!black] at ([xshift=-2.5mm,yshift=-2.2mm]n\i.north east) {$\checkmark$};
    \fi
    \ifnum\i=#1
      \draw[RedTitle, very thick, rounded corners=2pt]
        ([xshift=-0.6mm,yshift=-0.6mm]n\i.south west) rectangle ([xshift=0.6mm,yshift=0.6mm]n\i.north east);
      % keep the label inside the picture at both ends of the row
      \ifnum\i<3
        \node[font=\scriptsize\bfseries, text=RedTitle, anchor=south west] at ([yshift=1.2mm]n\i.north west)
          {$\blacktriangledown$\,you are here: #2};
      \else\ifnum\i>8
        \node[font=\scriptsize\bfseries, text=RedTitle, anchor=south east] at ([yshift=1.2mm]n\i.north east)
          {you are here: #2\,$\blacktriangledown$};
      \else
        \node[font=\scriptsize\bfseries, text=RedTitle, anchor=south] at ([yshift=1.2mm]n\i.north)
          {you are here: #2\,$\blacktriangledown$};
      \fi\fi
    \fi
  }
  % stage band under the sessions
  \foreach \a/\b/\lab/\c in {1/1/trunk/StageTrunk, 2/2/branches/StageBranch,
      3/5/the machine $\cdot$ classical methods/StageMachine, 6/6/the wall/StageWall,
      7/8/new tools/StageTools, 9/9/new data/StageData, 10/10/colleagues/StageAgents} {
    \fill[\c] ([yshift=-1.5mm]n\a.south west) rectangle ([yshift=-4.6mm]n\b.south east);
    \path ([yshift=-1.5mm]n\a.south west) -- ([yshift=-4.6mm]n\b.south east)
      node[midway, stage] {\lab};
  }
  \node[font=\footnotesize\itshape, text=black!70] at ([yshift=-9.5mm]$(n1.south)!0.5!(n10.south)$)
    {Build it on paper $\;\to\;$ solve it on the machine $\;\to\;$ push it to the frontier with AI};
\end{tikzpicture}
\end{center}
}


\graphicspath{{./figures/}}

\lstset{language=Python, basicstyle=\ttfamily\scriptsize, frame=none,
        backgroundcolor=\color{gray!8}, aboveskip=2pt, belowskip=2pt}

\decktitle{3}{A2}{Value Function Iteration: The First Complete Crossing}{Thursday, October 8, 2026}

\begin{document}

\makedecktitle

%=============================================================================
\begin{frame}{Where We Are}
\CourseMap{3}{3.A2}
\vspace{-1mm}
\begin{center}
\small \textbf{Block A:} \textbf{3.A1} what it means to solve a model $\;\cdot\;$ \textbf{3.A2} value function iteration on a grid $\;\cdot\;$ \textbf{3.A3} Euler equations and time iteration $\;\cdot\;$ \textbf{3.A4} the worked example, end to end.
\end{center}
\end{frame}

%=============================================================================
\begin{frame}[fragile]{Twenty Lines, and the Number They Have to Hit}
\begin{columns}[T]
\begin{column}{0.55\textwidth}
\begin{lstlisting}
import numpy as np
alpha, beta, tol = 0.36, 0.95, 1e-6
k = np.linspace(0.05, 0.60, 200)

C = (k**alpha)[:, None] - k[None, :]
U = np.full_like(C, -1e10)
np.log(C, out=U, where=C > 0)

V, n = np.zeros(200), 0
while True:
    M = U + beta * V[None, :]
    Vn, pol = M.max(axis=1), M.argmax(axis=1)
    d = np.max(np.abs(Vn - V))
    V, n = Vn, n + 1
    if d < tol:
        break

g = k[pol]
g_star = alpha * beta * k**alpha
print(n, np.max(np.abs(g / g_star - 1)))
\end{lstlisting}
\end{column}
\begin{column}{0.42\textwidth}
\small
Session 1 promised you this in 25 lines. It takes twenty. Every one of the six decisions of 3.A1 is in there, each made the crudest way available.
\medskip
\begin{examplebox}\footnotesize
\texttt{272 0.00988}
\end{examplebox}
\vspace{-1mm}
\begin{conceptbox}\footnotesize
272 sweeps, and a policy that is wrong by up to \textbf{1\%}.
\end{conceptbox}
\medskip
\footnotesize
The rest of this deck is two questions: \textbf{where did the 1\% come from}, and \textbf{why 272}?
\end{column}
\end{columns}
\end{frame}

%=============================================================================
\begin{frame}{What Those Lines Actually Do}
\begin{columns}[T]
\begin{column}{0.53\textwidth}
\small
\begin{enumerate}\footnotesize
  \item Build the grid $\{k_i\}_{i=1}^{N}$ and guess $v_0$.
  \item Build the \textbf{payoff matrix} once, outside the loop:
  \[ U_{ij}=u\big(zk_i^{\alpha}-k_j\big), \quad -\infty \text{ if } c\le0. \]
  \item Sweep: for every $i$,
  \[ v_{n+1}(k_i)=\max_{j}\Big\{U_{ij}+\beta\textstyle\sum_m \Pi_{jm} v_n(k_j,z_m)\Big\}. \]
  \item Stop when $\lVert v_{n+1}-v_n\rVert_\infty<\varepsilon$; read the policy off the \texttt{argmax}.
\end{enumerate}
\end{column}
\begin{column}{0.44\textwidth}
\begin{conceptbox}[title=\textbf{The one trick worth naming}]\small
$U$ does not depend on $v$, so it is built \textbf{once} and reused 272 times. Infeasible cells are set to $-10^{10}$ rather than tested inside the loop.
\end{conceptbox}
\medskip
\begin{cautionbox}\footnotesize
The $\max$ runs over \emph{grid points}, so the policy is forced onto the grid. That single decision is the source of most of the 1\%.
\end{cautionbox}
\end{column}
\end{columns}
\vspace{1mm}
\src{Algorithm and implementation steps adapted from QM Lec.~7, ``VFI: Algorithm Structure'' and ``VFI: Implementation Steps''.}
\end{frame}

%=============================================================================
\begin{frame}{Decision 2: Which Points, and How Many?}
\begin{columns}[T]
\begin{column}{0.50\textwidth}
\small
\textbf{Uniform.} $k_i=a+(i-1)\frac{b-a}{N-1}$. Simple, and interpolation error is uniform in $k$ --- which is the wrong thing to make uniform, because $v$ is far more curved at low $k$.
\medskip

\textbf{Logarithmic.} $k_i=a+(b-a)\dfrac{e^{\theta(i-1)/(N-1)}-1}{e^{\theta}-1}$, with $\theta\in[2,5]$: points bunched where the function bends.
\medskip

Both cost the same to evaluate. The second is almost always better, and almost always forgotten.
\end{column}
\begin{column}{0.47\textwidth}
\begin{cautionbox}[title=\textbf{Two boundary mistakes}]\footnotesize
\textbf{Too narrow} and the policy hits the top of the grid: $g(k_N)=k_N$ is a censored answer, not a solution.\\[1mm]
\textbf{Too wide} and most of your points sit where the model never visits, paying $O(N^2)$ for nothing.
\end{cautionbox}
\medskip
\begin{examplebox}\footnotesize
Always print $\max_i g(k_i)/k_N$. If it is 1, widen the grid before you believe anything.
\end{examplebox}
\end{column}
\end{columns}
\vspace{1mm}
\src{Grid families from QM Lec.~7, ``Grid Construction''. The curvature parameter is written $\theta$ here; the source calls it $\lambda$, which this course reserves for the invariant distribution.}
\end{frame}

%=============================================================================
\begin{frame}{Decision 3: Off the Grid, Straight Lines Only}
\begin{columns}[T]
\begin{column}{0.52\textwidth}
\small
Between two grid points,
\[
  \hat v(k)=v_i+\frac{v_{i+1}-v_i}{k_{i+1}-k_i}\,(k-k_i).
\]
\begin{itemize}\footnotesize
  \item Error is $O(h^2)$, with $h$ the spacing.
  \item It \textbf{preserves monotonicity and concavity} --- so a concave $v$ stays concave, and the maximization stays well behaved.
  \item Its derivative is discontinuous at every node. Anything that differentiates $\hat v$ will feel those kinks.
\end{itemize}
\end{column}
\begin{column}{0.45\textwidth}
\begin{conceptbox}[title=\textbf{Why not something smoother?}]\small
Cubic splines are $O(h^4)$ and $C^2$ --- and they \textbf{overshoot}, which can destroy concavity and send the maximizer somewhere absurd.
\end{conceptbox}
\medskip
\begin{takeawaybox}\footnotesize
Shape preservation beats formal accuracy inside a fixed-point loop. We stay linear today; the shape-preserving alternatives are built next week.
\end{takeawaybox}
\end{column}
\end{columns}
\vspace{1mm}
\src{Adapted from QM Lec.~7, ``Piecewise Linear Interpolation'' and ``Interpolation: Comparison''.}
\end{frame}

%=============================================================================
\begin{frame}{Decision 6: One Theorem, Used Two Ways}
\begin{columns}[T]
\begin{column}{0.51\textwidth}
\small
The contraction theorem gives \emph{two} inequalities, and they answer different questions.
\medskip

\textbf{How many iterations?} From $v_0$,
\[
  \lVert v_n-v^*\rVert\le\frac{\beta^{n}}{1-\beta}\,\lVert v_1-v_0\rVert
  \;\Rightarrow\; n\gtrsim\frac{\ln\varepsilon}{\ln\beta}.
\]

\textbf{How wrong am I now?} From the last step,
\[
  \lVert v_n-v^*\rVert\le\frac{\beta}{1-\beta}\,\lVert v_n-v_{n-1}\rVert .
\]
At $\beta=0.95$ that factor is \textbf{19}.
\end{column}
\begin{column}{0.46\textwidth}
\begin{center}\footnotesize
\begin{tabular}{@{}lrrr@{}}
\toprule
$\beta$ & $\ln\varepsilon/\ln\beta$ & measured & $\tfrac{\beta}{1-\beta}$ \\
\midrule
0.90 & 131.1 & 133 & 9 \\
0.95 & 269.3 & \textbf{272} & \textbf{19} \\
0.99 & 1374.6 & 1379 & 99 \\
\bottomrule
\end{tabular}
\end{center}
\vspace{1mm}
\begin{examplebox}\footnotesize
The count is a property of $\beta$ alone --- not of the grid, not of the model. Refining from $N=50$ to $N=800$ changed the iteration count by \textbf{zero}.
\end{examplebox}
\end{column}
\end{columns}
\vspace{1mm}
\src{$\beta^n/(1-\beta)$ from the contraction mapping theorem as stated in QM Lec.~2; the one-step form is 1.B2. Measured counts at $\varepsilon=10^{-6}$, $N=200$.}
\end{frame}

%=============================================================================
\begin{frame}{The Stopping Rule Is Not the Error}
\begin{columns}[c]
\begin{column}{0.53\textwidth}
\centering
\begin{tikzpicture}
\begin{axis}[
    width=7.9cm, height=5.3cm, xmin=0, xmax=280, ymin=-6.6, ymax=1.8,
    xlabel={\scriptsize iteration $n$}, ylabel={\scriptsize $\log_{10}$ error},
    tick label style={font=\scriptsize},
    axis x line*=bottom, axis y line*=left,
    legend style={font=\tiny, draw=none, fill=none, at={(0.98,0.97)}, anchor=north east}]
  \addplot[very thick, RedTitle] coordinates {(1,1.2994) (10,1.0953) (20,0.8725) (30,0.6498) (40,0.4270) (50,0.2042) (60,-0.0185) (70,-0.2413) (80,-0.4641) (90,-0.6868) (100,-0.9096) (110,-1.1323) (120,-1.3551) (130,-1.5779) (140,-1.8006) (150,-2.0234) (160,-2.2462) (170,-2.4689) (180,-2.6917) (190,-2.9145) (200,-3.1372) (210,-3.3600) (220,-3.5828) (230,-3.8055) (240,-4.0283) (250,-4.2510) (260,-4.4738) (270,-4.6966)};
  \addlegendentry{true error $\lVert v_n-v^h\rVert$}
  \addplot[very thick, orange!80!black, dashed] coordinates {(1,0.0925) (10,-0.1835) (20,-0.4062) (30,-0.6290) (40,-0.8518) (50,-1.0745) (60,-1.2973) (70,-1.5200) (80,-1.7428) (90,-1.9656) (100,-2.1883) (110,-2.4111) (120,-2.6339) (130,-2.8566) (140,-3.0794) (150,-3.3022) (160,-3.5249) (170,-3.7477) (180,-3.9705) (190,-4.1932) (200,-4.4160) (210,-4.6387) (220,-4.8615) (230,-5.0843) (240,-5.3070) (250,-5.5298) (260,-5.7526) (270,-5.9753)};
  \addlegendentry{what you monitor $\lVert v_n-v_{n-1}\rVert$}
  \draw[black!45, dotted, thick] (axis cs:0,-6) -- (axis cs:280,-6);
  \node[font=\tiny, text=black!55, anchor=south west] at (axis cs:4,-6) {$\varepsilon=10^{-6}$};
  \draw[<->, black!60, thick] (axis cs:150,-3.30) -- (axis cs:150,-2.02);
  \node[font=\tiny, anchor=west, text=black!60] at (axis cs:153,-2.66) {$\log_{10}19$};
\end{axis}
\end{tikzpicture}
\end{column}
\begin{column}{0.44\textwidth}
\small
The two lines are \textbf{parallel}, separated by exactly $\log_{10}19$. The quantity you watch is always a constant factor below the quantity you care about.
\medskip
\begin{center}\footnotesize
\begin{tabular}{@{}lr@{}}
\toprule
\multicolumn{2}{@{}l}{At $n=272$, when the loop stops:} \\
\midrule
$\lVert v_n-v_{n-1}\rVert$ & $9.6\times10^{-7}$ \\
guaranteed bound, $19\times$ & $1.81\times10^{-5}$ \\
\textbf{true error} & $\mathbf{1.81\times10^{-5}}$ \\
\bottomrule
\end{tabular}
\end{center}
\vspace{1mm}
\begin{cautionbox}\footnotesize
The $19\times$ bound is not pessimistic here --- it is \textbf{attained}. Ask for $10^{-6}$ and you get $1.8\times10^{-5}$.
\end{cautionbox}
\end{column}
\end{columns}
\end{frame}

%=============================================================================
\begin{frame}{And You Can Do Much Better, for Free}
\begin{columns}[T]
\begin{column}{0.54\textwidth}
\small
Let $d_{\min}=\min_i\big(v_{n+1}-v_n\big)(k_i)$ and $d_{\max}$ the maximum. Then \emph{both} sides are bounded:
\[
  v_{n+1}+\tfrac{\beta}{1-\beta}d_{\min}\;\le\;v^*\;\le\;v_{n+1}+\tfrac{\beta}{1-\beta}d_{\max}.
\]
The naive rule throws away $d_{\min}$ and uses $\lVert\cdot\rVert_\infty$ for both. Keeping both ends costs two extra reductions per sweep and gives a \textbf{bracket} instead of a bound --- take its midpoint and you have a free accuracy upgrade.
\medskip

At $n=272$ that bracket is $2.7\times10^{-13}$ wide, against a true error of $1.8\times10^{-5}$.
\end{column}
\begin{column}{0.43\textwidth}
\begin{cautionbox}[title=\textbf{Read that honestly}]\footnotesize
Eight orders of magnitude is \emph{this model} flattering us: with $\delta=1$ and log utility the increment $v_{n+1}-v_n$ is nearly constant in $k$, so $d_{\max}\approx d_{\min}$.
\end{cautionbox}
\medskip
\begin{takeawaybox}\footnotesize
In general the bracket is narrower, not necessarily \emph{that} much narrower. But it is never worse than the one-sided bound, and it is two lines of code.
\end{takeawaybox}
\end{column}
\end{columns}
\vspace{1mm}
\src{MacQueen (1966); Porteus (1971). Statement adapted from QM Lec.~8, ``MacQueen--Porteus Bounds''.}
\end{frame}

%=============================================================================
\begin{frame}{What It Costs}
\begin{columns}[T]
\begin{column}{0.50\textwidth}
\small
Each sweep compares every grid point against every candidate choice: $O(N^2)$ per iteration, $O(N^2 N_z^2)$ with a shock. The iteration count is $O(\ln\varepsilon/\ln\beta)$ and does not depend on $N$. So the total is
\[
  \underbrace{N^2}_{\text{grid}}\times\underbrace{\frac{\ln\varepsilon}{\ln\beta}}_{\text{patience}} .
\]
\begin{center}\footnotesize
\begin{tabular}{@{}lrr@{}}
\toprule
$N$ & sweeps & comparisons \\
\midrule
200 & 272 & 10.9 million \\
400 & 272 & 43.5 million \\
800 & 272 & 174 million \\
\bottomrule
\end{tabular}
\end{center}
\end{column}
\begin{column}{0.47\textwidth}
\begin{cautionbox}[title=\textbf{Two ways this becomes fatal}]\footnotesize
\textbf{Patience.} At a quarterly $\beta=0.99$ the sweep count is 1{,}379 instead of 272 --- a $5\times$ bill for changing the calendar.\\[1mm]
\textbf{Dimension.} A second continuous state makes the grid $N^2$ and the sweep $N^4$. 200 points per dimension is $1.6\times10^{9}$ comparisons \emph{per sweep}.
\end{cautionbox}
\medskip
\begin{examplebox}\footnotesize
Nothing here is wasted motion --- it is the crude choices being paid for. Every one has a cheaper replacement, and that is next week.
\end{examplebox}
\end{column}
\end{columns}
\vspace{1mm}
\src{Cost accounting adapted from QM Lec.~8, ``Recap: Standard VFI Algorithm'' and ``Computational Bottlenecks''.}
\end{frame}

%=============================================================================
\begin{frame}{Refining the Grid --- Until It Stops Helping}
\begin{columns}[c]
\begin{column}{0.51\textwidth}
\centering
\begin{tikzpicture}
\begin{axis}[
    width=7.6cm, height=5.3cm, xmin=-3.30, xmax=-1.80, ymin=-5.6, ymax=-1.45,
    xlabel={\scriptsize $\log_{10} h$}, ylabel={\scriptsize $\log_{10}$ max error},
    tick label style={font=\scriptsize},
    axis x line*=bottom, axis y line*=left,
    legend style={font=\tiny, draw=none, fill=none, at={(0.02,0.97)}, anchor=north west}]
  \addplot[very thick, RedTitle, mark=*, mark size=1.4pt] coordinates {(-1.9498,-2.9500) (-2.2553,-3.2729) (-2.5585,-3.7887) (-2.8606,-4.5513) (-3.1622,-4.7463)};
  \addlegendentry{value, $\varepsilon=10^{-6}$}
  \addplot[very thick, orange!80!black, mark=square*, mark size=1.2pt] coordinates {(-1.9498,-2.9430) (-2.2553,-3.2584) (-2.5585,-3.7427) (-2.8606,-4.3349) (-3.1622,-5.2736)};
  \addlegendentry{value, $\varepsilon=10^{-11}$}
  \addplot[very thick, black!55, mark=triangle*, mark size=1.4pt] coordinates {(-1.9498,-2.2168) (-2.2553,-2.4953) (-2.5585,-2.7599) (-2.8606,-3.0396) (-3.1622,-3.3499)};
  \addlegendentry{policy}
  \draw[black!45, dotted, thick] (axis cs:-3.30,-4.7212) -- (axis cs:-1.80,-4.7212);
  \node[font=\tiny, text=black!55, anchor=north east] at (axis cs:-1.84,-4.76) {truncation floor $19\varepsilon$};
\end{axis}
\end{tikzpicture}
\end{column}
\begin{column}{0.46\textwidth}
\small
Santos \& Vigo-Aguiar (1998): the value error is $O(h^2)$ and the policy error $O(h)$. Measured slopes at $\varepsilon=10^{-11}$: \textbf{1.89} and \textbf{0.93}.
\medskip

At $\varepsilon=10^{-6}$ the value curve \textbf{flattens}. Going from $N=400$ to $N=800$ --- four times the work --- moves the value error from $2.8\times10^{-5}$ to $1.8\times10^{-5}$, because it has hit the $19\varepsilon$ floor from 3.A1.
\medskip
\begin{takeawaybox}\footnotesize
The budgets compose. Refining a grid while the tolerance is loose buys nothing, and you would never see it without running both.
\end{takeawaybox}
\end{column}
\end{columns}
\end{frame}

%=============================================================================
\begin{frame}{Graded Against the Closed Form}
\begin{columns}[c]
\begin{column}{0.50\textwidth}
\centering
\begin{tikzpicture}
\begin{axis}[
    width=7.4cm, height=5.2cm, xmin=0.213, xmax=0.315, ymin=0.194, ymax=0.227,
    xlabel={\scriptsize capital $k$ (36 consecutive grid points)},
    ylabel={\scriptsize policy $k'$},
    tick label style={font=\scriptsize}, scaled ticks=false,
    x tick label style={/pgf/number format/fixed, /pgf/number format/precision=2},
    y tick label style={/pgf/number format/fixed, /pgf/number format/precision=2},
    axis x line*=bottom, axis y line*=left,
    legend style={font=\tiny, draw=none, fill=none, at={(0.98,0.03)}, anchor=south east}]
  \addplot[const plot, thick, RedTitle] coordinates {(0.21583,0.19648) (0.21859,0.19925) (0.22136,0.19925) (0.22412,0.19925) (0.22688,0.19925) (0.22965,0.20201) (0.23241,0.20201) (0.23518,0.20201) (0.23794,0.20477) (0.24070,0.20477) (0.24347,0.20477) (0.24623,0.20754) (0.24899,0.20754) (0.25176,0.20754) (0.25452,0.21030) (0.25729,0.21030) (0.26005,0.21030) (0.26281,0.21307) (0.26558,0.21307) (0.26834,0.21307) (0.27111,0.21307) (0.27387,0.21583) (0.27663,0.21583) (0.27940,0.21583) (0.28216,0.21583) (0.28492,0.21859) (0.28769,0.21859) (0.29045,0.21859) (0.29322,0.22136) (0.29598,0.22136) (0.29874,0.22136) (0.30151,0.22136) (0.30427,0.22412) (0.30704,0.22412) (0.30980,0.22412) (0.31256,0.22412)};
  \addlegendentry{VFI on the grid}
  \addplot[thick, black!55, dashed] coordinates {(0.21583,0.19693) (0.21859,0.19783) (0.22136,0.19873) (0.22412,0.19962) (0.22688,0.20050) (0.22965,0.20138) (0.23241,0.20225) (0.23518,0.20311) (0.23794,0.20396) (0.24070,0.20481) (0.24347,0.20566) (0.24623,0.20650) (0.24899,0.20733) (0.25176,0.20815) (0.25452,0.20897) (0.25729,0.20979) (0.26005,0.21059) (0.26281,0.21140) (0.26558,0.21220) (0.26834,0.21299) (0.27111,0.21377) (0.27387,0.21456) (0.27663,0.21533) (0.27940,0.21611) (0.28216,0.21687) (0.28492,0.21764) (0.28769,0.21839) (0.29045,0.21915) (0.29322,0.21989) (0.29598,0.22064) (0.29874,0.22138) (0.30151,0.22211) (0.30427,0.22284) (0.30704,0.22357) (0.30980,0.22429) (0.31256,0.22501)};
  \addlegendentry{$\alpha\beta k^{\alpha}$}
\end{axis}
\end{tikzpicture}
\end{column}
\begin{column}{0.47\textwidth}
\small
Zoom in and the agreement becomes a \textbf{staircase}. Across these 36 states the computed policy takes only \textbf{11 distinct values} --- it can return nothing but a grid point.
\medskip
\begin{center}\footnotesize
\begin{tabular}{@{}lr@{}}
\toprule
max $|v-v^*|$ & $1.6\times10^{-4}$ \\
max relative policy error & $0.99\%$ \\
grid spacing as a share of $k'$ & $\approx1.4\%$ \\
\bottomrule
\end{tabular}
\end{center}
\medskip
\begin{cautionbox}\footnotesize
The error is a \textbf{sawtooth}, not noise, and it does not shrink with more iterations. It is decision 4 --- searching the grid --- and no amount of patience will fix it.
\end{cautionbox}
\end{column}
\end{columns}
\end{frame}

%=============================================================================
\begin{frame}{Decision 5: Where Does the Expectation Come From?}
\begin{columns}[T]
\begin{column}{0.52\textwidth}
\small
Add productivity $z\in\{z_1,\dots,z_{N_z}\}$ with transition matrix $\Pi$. The sweep becomes
\[
  v_{n+1}(k_i,z_j)=\max_{l}\Big\{U_{ijl}+\beta\sum_{m}\Pi_{jm}v_n(k_l,z_m)\Big\},
\]
and the expectation is a matrix product: \texttt{EV = V @ Pi.T}, computed once per sweep rather than inside the maximization.
\medskip

With $z\in\{0.95,1.05\}$ and $\Pi$ symmetric with persistence $0.9$, the closed form still holds state by state --- so we can still grade it.
\end{column}
\begin{column}{0.45\textwidth}
\begin{cautionbox}[title=\textbf{But where did $\Pi$ come from?}]\footnotesize
We \emph{wrote it down}. Real work starts from an estimated AR(1),
$\ln z'=\rho\ln z+\sigma\epsilon'$, and must turn it into a chain that matches the persistence and the variance.
\end{cautionbox}
\medskip
\begin{examplebox}\footnotesize
That is a genuine gap, not a detail: the two standard constructions disagree badly when $\rho>0.95$, which is exactly the range macro uses. We build both next week.
\end{examplebox}
\end{column}
\end{columns}
\vspace{1mm}
\src{Stochastic sweep adapted from QM Lec.~7, ``Time Iteration with Stochastic Shocks'' and the Lec.~8 recap; the chain used here is the one introduced in 1.B1.}
\end{frame}

%=============================================================================
\begin{frame}{Six Crude Choices, Six Gaps}
\renewcommand{\arraystretch}{1.15}
\begin{center}\footnotesize
\begin{tabularx}{\linewidth}{@{}>{\raggedright\arraybackslash}p{2.7cm} >{\raggedright\arraybackslash}X >{\raggedright\arraybackslash}p{3.2cm} c@{}}
\toprule
\textbf{What we did} & \textbf{What it cost} & \textbf{What we need} & \\
\midrule
Searched the grid for the $\max$ & the 1\% sawtooth; $O(N^2)$ per sweep & optimization on an interval & S4 \\
Linear interpolation & $O(h^2)$, kinked derivatives & approximation theory & S4 \\
Wrote $\Pi$ down by hand & no link to any estimated process & discretizing an AR(1) & S4 \\
Naive stopping rule & $19\times$ the error you asked for & two-sided bounds & today \\
$N^2\times272$ comparisons & 174 million at $N=800$ & acceleration, better operators & S4 \\
One state variable & nothing yet --- and that is the problem & perturbation, projection & S5 \\
\bottomrule
\end{tabularx}
\end{center}
\vspace{1mm}
\begin{takeawaybox}\small
Every row is a numerical technique that exists \textbf{because} a crude choice hurt. That is the honest reason the next two sessions look like a mathematics course: each tool on the list was invented to pay off one of these lines.
\end{takeawaybox}
\end{frame}

%=============================================================================
\begin{frame}{Takeaways}
\begin{takeawaybox}
\begin{itemize}\small
  \item Twenty lines take you from the Bellman equation to a policy function. The crossing is \textbf{not} the hard part --- knowing how wrong you are is.
  \item The iteration count belongs to $\beta$ alone: 272 sweeps at $\beta=0.95$, 1{,}379 at $0.99$, whatever the grid. The \textbf{accuracy} belongs to the grid.
  \item The quantity you monitor, $\lVert v_n-v_{n-1}\rVert$, sits a factor $\beta/(1-\beta)$ below the error you care about --- $19\times$ here, and that bound was \textbf{attained}, not merely respected.
  \item Error budgets compose. At $\varepsilon=10^{-6}$, quadrupling the grid bought nothing, because the tolerance had already put a floor under the value error.
  \item The 1\% policy error is a \textbf{sawtooth}, not noise: $\arg\max$ can only return a grid point. Diagnose the shape of an error, not just its size.
\end{itemize}
\end{takeawaybox}
\end{frame}

%=============================================================================
\begin{frame}{Bridge: The Same Model, Down the Other Road}
\begin{columns}[T]
\begin{column}{0.53\textwidth}
\small
Value function iteration never used the first-order condition. It maximized, brutally, at every point of every sweep --- and the price was a policy quantized to the grid.
\medskip

But we know the optimum satisfies
\[
  u'(c)=\beta\,\mathbb{E}\big[u'(c')R(k',z')\big].
\]
That is an equation the policy must satisfy \emph{exactly}, and solving an equation is a different --- often cheaper, often sharper --- numerical problem than searching for a maximum.
\end{column}
\begin{column}{0.44\textwidth}
\begin{conceptbox}[title=\textbf{Next (3.A3)}]\small
Iterate on the \textbf{policy} instead of the value: the Coleman and Greenwood--Huffman operators, one root-find per grid point.\\[1mm]
Then the comparison that matters --- two methods, one model, measured against the same closed form --- and the list of what still breaks when the model gets harder.
\end{conceptbox}
\end{column}
\end{columns}
\end{frame}

\end{document}
